% Generated draft: requires mathematical and source review.
\documentclass[11pt]{article}
\usepackage[margin=25mm]{geometry}
\usepackage{fontspec}
\setmainfont{Latin Modern Roman}
\usepackage{amsmath,amssymb,mathtools,unicode-math}
\setmathfont{Latin Modern Math}
\usepackage{enumitem,needspace,xurl,hyperref}
\setlength{\parindent}{0pt}
\setlength{\parskip}{5pt}
\setlength{\emergencystretch}{3em}
\newcommand{\N}{\mathbb{N}}
\newcommand{\Z}{\mathbb{Z}}
\newcommand{\Q}{\mathbb{Q}}
\newcommand{\R}{\mathbb{R}}
\newcommand{\C}{\mathbb{C}}
\newcommand{\F}{\mathbb{F}}
\newcommand{\A}{\mathbb{A}}
\newcommand{\PP}{\mathbb{P}}
\newcommand{\E}{\mathbb{E}}
\newcommand{\eps}{\varepsilon}
\newcommand{\dd}{\,\mathrm{d}}
\begin{document}
\section*{UnsolvedMath ID: 30006279}
\textbf{Source problem code:} catalogue entry
\subsection{Definitions and mathematical statement}
The supplied record asks for a universal-cover decomposition of a Kawamata log terminal (klt) pair with nef anticanonical class. Its background discusses projective pairs and a Beauville--Bogomolov-type decomposition, but the extracted passage does not specify enough hypotheses or the proposed factors to formulate a complete conjecture.

For terminology, in the usual complex projective setting a pair \((X,\Delta)\) consists of a normal projective complex variety \(X\) and an effective divisor \(\Delta\), with coefficients and Cartier conditions appropriate to the chosen category. For example, for an effective \(\R\)-divisor \(\Delta\) with \(K_X+\Delta\) \(\R\)-Cartier, the pair is klt if every divisorial discrepancy satisfies
\[
a(E,X,\Delta)>-1
\qquad\text{for every prime divisor }E\text{ over }X.
\]
An \(\R\)-Cartier divisor \(D\) on a projective variety is nef if
\[
D\cdot C\geq 0
\qquad\text{for every integral projective curve }C\subseteq X.
\]
These definitions explain standard terminology; the complex ground field, divisor coefficients, and precise Cartier assumptions are not explicitly fixed by the extracted question.

\textbf{Source question, with unresolved qualifications.} Describe the decomposition of the universal cover associated with a klt pair having nef anticanonical class, in the Beauville--Bogomolov-type sense mentioned in the supplied background. The original passage directs the reader to \([\mathrm{MW21},\text{ Conjecture }1.5]\).

For a pair, the anti-log-canonical divisor is \(-(K_X+\Delta)\), whereas the anticanonical divisor of the underlying variety is \(-K_X\). The supplied background uses the former, while the statement says only ``anticanonical class''; these conditions cannot be silently identified. Likewise, if \(X\) is complex, the topological universal cover of \(X^{\mathrm{an}}\) and that of its smooth locus \(X_{\mathrm{reg}}^{\mathrm{an}}\) are distinct possible targets, wherever the relevant covering exists. The source does not select between them or another pair-associated covering. It also does not specify the proposed product factors, their properties, the treatment of the boundary, or the category of decomposition. Consequently, no precise product formula is asserted here; the cited conjecture must be recovered and its hypotheses checked before completing the formulation.
\subsection{Short English statement}
The supplied question asks for a Beauville--Bogomolov-type decomposition of the universal cover associated with a klt pair under a nefness assumption. The divisor required to be nef, the intended cover, and the proposed factors remain unspecified. The supplied record labels the problem open; this draft does not independently verify that status.
\subsection{Sources}
\begin{itemize}
\item Supplied local UnsolvedMath registry record, permanent ID \(30006279\), source problem code \texttt{catalogue entry}, titled ``Universal-Cover Decomposition of Klt Pairs with Nef Anticanonical Class.'' The record belongs to the Oberwolfach Reports Open Problems collection, assigns the year 2025, and names no individual proposer.
\item \emph{Recent Trends in Algebraic Geometry} (2025), Oberwolfach Reports, DOI 10.4171/owr/2025/28, pp.\ 1451--1516. Supplied report link and literature-source label ``original OWR report'': \url{https://doi.org/10.4171/owr/2025/28}. The supplied passage refers to \([\mathrm{MW21},\text{ Conjecture }1.5]\), but its bibliography entry and conjecture text are not included.
\item Supplied literature source: ``Matsumura--Wang, non-vanishing and structure,'' \url{https://doi.org/10.1017/nmj.2024.21}. The record does not establish the precise bibliographic relationship between this link and the label \([\mathrm{MW21}]\).
\item Literature assessment dated \texttt{2026-08-22}: the supplied curation labels the problem open and reports that structure theorems for projective klt pairs with nef anti-log-canonical divisor give locally constant fibrations and reduce part of the problem, while the requested universal-cover decomposition remains unresolved in the stated generality. This is supplied dated curation, not an independently verified literature assessment.
\item The registry assigns the label \texttt{corrected\_verified} and states that the official report was checked, the question isolated, and its notation repaired. These are supplied source-verification claims; no independent source consultation or mathematical verification is asserted in this draft.
\end{itemize}
\end{document}
